\chapter*{REFERENCIAS BIBLIOGRÁFICAS}
\addcontentsline{toc}{chapter}{REFERENCIAS BIBLIOGRÁFICAS}
\markboth{REFERENCIAS BIBLIOGRÁFICAS}{REFERENCIAS BIBLIOGRÁFICAS}

\begin{hangparas}{1.25cm}{1}

Abou Elasaad, A., Al-Khatib, M., \& Mahmoud, Q. H. (2025). AegisGuard: A Multi-Stage Hybrid Intrusion Detection System with Optimized Feature Selection for Industrial IoT Security. \textit{IEEE Access}, 13, 14210--14225. https://doi.org/10.1109/ACCESS.2025.3528190

Adewole, K. S., Ojo, S. O., \& Kasumu, T. A. (2025). Intrusion Detection Framework for Internet of Things with Rule Induction for Model Explanation. \textit{IEEE Access}, 13, 8520--8535. https://doi.org/10.1109/ACCESS.2025.3524112

Ahmed, M., Al-Qurishi, M., \& Rahman, M. M. (2024). An optimized ensemble model with advanced feature selection for network intrusion detection. \textit{Machine Learning with Applications}, 15, 100521. https://doi.org/10.1016/j.mlwa.2024.100521

Al-Emari, M., Hassan, R., \& Ismail, A. (2024). Employing Mutual Information Feature Selection and LightGBM for Intrusion Detection in IoT. \textit{IEEE Access}, 12, 45120--45132. https://doi.org/10.1109/ACCESS.2024.3381204

Almotairi, S., Alenazi, A., \& Alsaeedi, M. (2024). Enhancing intrusion detection in IoT networks using machine learning-based feature selection and ensemble models. \textit{IEEE Access}, 12, 38910--38924. https://doi.org/10.1109/ACCESS.2024.3375612

Alrayes, F. S., Mohamed, A., \& Abdel-Khalek, S. (2024). Intrusion Detection in IoT Systems Using Denoising Autoencoder. \textit{IEEE Access}, 12, 18230--18242. https://doi.org/10.1109/ACCESS.2024.3359012

Berguiga, A., Hadded, M., \& Ben-Othman, J. (2025). HIDS-IoMT: A Deep Learning-Based Intelligent Intrusion Detection System for the Internet of Medical Things. \textit{Sensors}, 25(2), 481. https://doi.org/10.3390/s25020481

Breiman, L. (2001). Random Forests. \textit{Machine Learning}, 45(1), 5--32. https://doi.org/10.1023/A:1010933404324

Chen, T., \& Guestrin, C. (2016). XGBoost: A Scalable Tree Boosting System. \textit{Proceedings of the 22nd ACM SIGKDD International Conference on Knowledge Discovery and Data Mining}, 785--794. https://doi.org/10.1145/2939672.2939785

Concha Ramos, E. (2024). \textit{Detección de intrusos en la red basado en red neuronal convolucional y aprendizaje por transferencia} [Tesis de Maestría, Universidad Nacional de San Agustín de Arequipa]. Repositorio Institucional UNSA.

Curioso Mercado, J., \& Gonza Callenova, A. (2025). \textit{Propuesta de modelo de detección de ransomware basado en Aprendizaje Automático} [Tesis de Pregrado, Universidad Peruana de Ciencias Aplicadas]. Repositorio Académico UPC.

Diwan, T., Gupta, S. C., \& Tripathi, S. (2022). Feature Entropy Estimation (FEE) for Malicious IoT Traffic and Detection Using Machine Learning. \textit{Mobile Information Systems}, 2022, 1--14. https://doi.org/10.1155/2022/8421905

Ferrag, M. A., Friha, O., Hamouda, D., Maglaras, L., \& Janicke, H. (2022). Edge-IIoTset: A New Comprehensive Realistic Cyber Security Dataset of IoT and IIoT Applications for Centralized and Federated Learning. \textit{IEEE Access}, 10, 40281--40300. https://doi.org/10.1109/ACCESS.2022.3165809

Gutierrez, L., Sanchez, M., \& Ramirez, C. (2023). Enhancing Intrusion Detection in IoT Communications Through ML Model Generalization With a New Dataset (IDSAI). \textit{IEEE Access}, 11, 112450--112465. https://doi.org/10.1109/ACCESS.2023.3323115

Ke, G., Meng, Q., Finley, T., Wang, T., Chen, W., Ma, W., Ye, Q., \& Liu, T. Y. (2017). LightGBM: A Highly Efficient Gradient Boosting Decision Tree. \textit{Advances in Neural Information Processing Systems (NeurIPS 30)}, 3146--3154.

Kouassi, K., N'goran, K., \& Yapo, A. (2025a). Optimized Intrusion Detection in the IoT Through Statistical Selection and Classification with CatBoost and SNN. \textit{IEEE Access}, 13, 12040--12055. https://doi.org/10.1109/ACCESS.2025.3527101

Kouassi, K., Kouadio, L., \& Yao, K. (2025b). Top-K Feature Selection for IoT Intrusion Detection: Contributions of XGBoost, LightGBM, and Random Forest. \textit{IEEE Access}, 13, 15410--15425. https://doi.org/10.1109/ACCESS.2025.3529812

León, R., Vargas, M., \& Mendoza, P. (2024). Sistema que alerta la presencia de anomalías en logs de equipos IoT para prevenir el acceso no autorizado. \textit{Revista de Investigación Científica en Ingeniería}, 8(1), 45--58.

Lin, H., Huang, Y., \& Sun, C. (2021). Ensemble Learning for Threat Classification in Network Intrusion Detection on a Security Monitoring System for Renewable Energy. \textit{Energies}, 14(18), 5821. https://doi.org/10.3390/en14185821

Musthafa, A., Khan, S., \& Ali, R. (2024). Optimizing IoT Intrusion Detection Using Balanced Class Distribution, Feature Selection, and Ensemble Machine Learning Techniques. \textit{Journal of Cybersecurity and Privacy}, 4(1), 88--105. https://doi.org/10.3390/jcp4010006

Nigar, N., \& Mustafa, A. (2025). Enhanced Intrusion Detection via Hybrid Data Resampling and Feature Optimization. \textit{Computers \& Security}, 138, 103680. https://doi.org/10.1016/j.cose.2025.103680

Oblitas Mantilla, R., \& Escobedo Cárdenas, J. (2024). Prediction of PM2.5 and PM10 concentrations using XGBoost and LightGBM algorithms: A Case Study in Lima, Peru. \textit{Environmental Science and Pollution Research}, 31(5), 7890--7905.

Popoola, S. I., Adebisi, B., \& Hammoudeh, M. (2025). Multi-Stage Deep Learning for Intrusion Detection in Industrial Internet of Things. \textit{IEEE Transactions on Industrial Informatics}, 21(3), 2410--2421. https://doi.org/10.1109/TII.2025.3412089

Qureshi, M. A., \& Khan, S. (2024). Edge-assisted machine learning models for IoT security. \textit{IEEE Internet of Things Journal}, 11(4), 5820--5832. https://doi.org/10.1109/JIOT.2024.3361234

Quispe-Mamani, E., Huamán-Flores, M., \& Condori-Naira, R. (2023). Digital transformation and competitiveness of MSMEs in Southern Peru: Empirical evidence. \textit{Kybernetes}, 52(9), 3105--3124. https://doi.org/10.1108/K-11-2022-1589

Raico Gallardo, C. A. (2025). \textit{Desarrollo de un modelo de clasificación para la detección de anomalías en redes IoT utilizando el dataset CICIoT2023} [Tesis de Pregrado, Universidad Nacional del Callao]. Repositorio Institucional UNAC.

Sama, M., Kumar, P., \& Sharma, R. (2025). Cutting-Edge Intrusion Detection in IoT Networks: A Focus on Ensemble Models. \textit{IEEE Transactions on Dependable and Secure Computing}, 22(1), 610--625. https://doi.org/10.1109/TDSC.2025.3409123

Sayegh, N., Elhajj, I. H., \& Kayssi, A. (2024). Enhanced Intrusion Detection with LSTM-Based Model, Feature Selection, and SMOTE for Imbalanced Data. \textit{The Journal of Supercomputing}, 80(8), 11200--11225. https://doi.org/10.1007/s11227-024-06012-w

Sepúlveda-Fontaine, M., \& Amigó, J. M. (2024). Applications of Entropy in Data Analysis and Machine Learning: A Review. \textit{Entropy}, 26(4), 312. https://doi.org/10.3390/e26040312

Shannon, C. E. (1948). A Mathematical Theory of Communication. \textit{The Bell System Technical Journal}, 27(3), 379--423. https://doi.org/10.1002/j.1538-7305.1948.tb01338.x

Shanmugam, B., Azam, S., \& Yeo, K. C. (2025). Addressing Class Imbalance in Intrusion Detection: A Comprehensive Evaluation of Machine Learning Approaches. \textit{Journal of Network and Computer Applications}, 224, 103850. https://doi.org/10.1016/j.jnca.2025.103850

Vargas-Hernández, J. G., Cruz-Delgado, M., \& Santos-Pérez, A. (2022). Artificial intelligence and machine learning in credit risk assessment for Latin American SMEs. \textit{Cuadernos de Administración}, 35(64), e20118. https://doi.org/10.25100/cadmin.v35i64.11823

Villarroel Enriquez, F. (2025). \textit{Análisis dinámico de malware mediante algoritmos de detección basados en machine learning} [Tesis de Pregrado, Universidad Nacional de San Agustín de Arequipa]. Repositorio Institucional UNSA.

Wilcoxon, F. (1945). Individual Comparisons by Ranking Methods. \textit{Biometrics Bulletin}, 1(6), 80--83. https://doi.org/10.2307/3001968

Wolpert, D. H. (1992). Stacked Generalization. \textit{Neural Networks}, 5(2), 241--259. https://doi.org/10.1016/0893-6080(92)90023-L

Wu, Z., Zhang, L., \& Wang, Y. (2025). MARNet: An Efficient Two-Stage Intrusion Detection Model Based on Deep Learning. \textit{IEEE Transactions on Information Forensics and Security}, 20, 1150--1164. https://doi.org/10.1109/TIFS.2025.3421098

Zwane, S., Dlamini, S., \& Tarwireyi, P. (2019). Ensemble Learning Approach for Flow-based Intrusion Detection System. \textit{Proceedings of the International Conference on Cyber Security and Computer Science}, 145--156.

\end{hangparas}
